\section*{РЕЗЮМЕ СТАРТАП-ПРОЕКТА}
\addcontentsline{toc}{section}{РЕЗЮМЕ СТАРТАП-ПРОЕКТА}
Название: «\Тема\ \ТемаВтораяСтрока». Проект представляет собой систему, предназначенную  для автоматизации и оптимизации бизнес-процессов заведений общественного питания. Она цифровизирует такие задачи как прием, обработка и оплата заказа; хранение данных о сотрудниках; оформление графика работы персонала; создание и редактирование технологических карт, ведение бухгалтерского учёта. Были проанализированы современные аналоги и выявлены их \textit{отрицательные стороны}:
\begin{enumerate}
	\item Сложность настройки и использования. Большинство программно-информационных систем для оптимизации ресторанного бизнеса наделены огромный количеством функций, которые требуют настройки под каждое заведение персонально. Данный процесс может занять от нескольких часов до нескольких дней.
	\item Высокая стоимость: самые распространённые системы подключаются по ежемесячной, еженедельной, годовой подписках. Каждое устройство и терминал оплачивается отдельно, вследствие чего для доступа к максимальному количеству функций, требуется затратить большое количество средств на подписку. 
	\item Многие программно-информационные системы были созданы в промежуток с 1992 года по 2006 год. И с тех пор интерфейс этих программ практически не изменился. В них преобладает серый цвет, хаотично располагаются кнопки, что является причиной для путаницы персонала и требует больших временных затрат на обучение.
	\item В некоторых системах возможность просмотра заказов предыдущих дней отсутствует или же является временно-затратным процессом.
	\item Возможность не только разделить оплаты по типам, но и использовать один тип оплат заказа несколько раз. Это позволяет высчитать сдачу для каждого гостя персонально, при раздельном счёте.
\end{enumerate} 


Созданная программно-информационная система позволит избежать вышеперечисленных трудностей без негативного влияния на процесс работы и скорость выполнения поставленных задач.

\textit{Актуальность} разрабатываемой системы связана с потребностью заведений общественного питания в автоматизации бизнес-процессов, оптимизации работы с заказом, эффективном управлении данными предприятия, сотрудников и позициями меню. 

\textit{Целью} является разработка программно-информационной системы для оптимизации процессов в сфере ресторанного бизнеса. 


\textit{Отличительные особенности} программно-информационной системы по сравнению с уже существующими решениями и программами:
\begin{enumerate}
	\item Простой и интуитивно понятный графический интерфейс с модулями, разделенными между должностями сотрудников.
	\item Возможность создания графиков работы персонала и сохранение их в формате .xlxs.
	\item Процесс создания технологической карты разделен на блоки: добавление категории, названия блюда, ингредиентов.
	\item Управление заказами с возможностью их просмотра за предыдущие смены.
\end{enumerate}


В разработанной программно-информационной системе управления ресторанным бизнесом реализованы такие задачи как: просмотр, создание и изменение данных о персонале, блюдах, столах, типах оплаты, возможность составлять графики для сотрудников, формировать финансовые отчёты, создавать и редактировать заказ.

Идея создание проекта появилась в сентябре 2023 года, разработка была начата в сентябре 2024 года. А с марта 2025 производилось тестирование системы путем формирования заказов. Проект планируется закончить к началу 2026 года. Также будет осуществлена регистрация программно-информационной системы для избежания плагиата и обеспечения исключительности прав.
\newpage
\textit{Бизнес-цели:}
 \begin{enumerate}
 	\item Добавить возможность работы с клиентской базой. То есть регистрацию гостей и бонусную систему.
 	\item Связать программный продукт с сайтом qr-меню и бронированием столов.
 	\item Создать мобильное приложение для работы официантов и хостес.
 	\item Добавить функциональную возможность связи программного продукта с социальными сетями компании с целью публикации актуальной информации об акциях, скидках заведения общественного питания.
 	\item Внедрить готовую систему в ресторан.
 \end{enumerate}




